\subsection{Symbolic Axis Mapping}
\label{sec:axis-mapping}
Expressing operations through shared mathematical primitives provides a
common basis for reasoning about tensor rewrites. TEPL composes their
value expressions and normalizes the LHS and RHS under the permitted
arithmetic laws to check equivalence.

Consider \code{T=dot(X,Y); O=add(T,Z)}, with batch group $\vec B$ and no broadcasting.
The dot's value expression is \code{Reduce(K,Product(X,Y))}, and addition's
is \code{Add(T,Z)}. Substituting the dot expression for \code{T} gives
\code{Add(Reduce(K,Product(X,Y)),Z)}.
This describes the arithmetic, but the operations initially use different
axis groups: the dot yields $\vec B\catseq\vec I\catseq\vec J$, while addition binds
all axes of $T$ and $Z$ to a single group $\vec S$ and yields $\vec S$ as its output
axis order. Thus $\vec S$ contains all axes of $T$, $Z$, and $O$ in their shared
order. The compiler must identify $\vec S=\vec B\catseq\vec I\catseq\vec J$ so that addition
reads the dot result and $Z$ at corresponding coordinates.

Symbolic axis mapping establishes this correspondence. Once the bindings
are aligned, the composed expression describes, for batched matrices,
\begin{equation}
 O_{bij}=\sum_k X_{bik}Y_{bkj}+Z_{bij}.
\end{equation}
Here the operand names in the value expression stand for scalar reads at
the aligned coordinates. This lets TEPL compare composed computations
using both their arithmetic and their axis relationships.
The following subsections explain how these mappings are built and composed
when group lengths and positions are symbolic.

\subsubsection{Building Mappings Across Operations}
\label{sec:mapping-tables}
\sidefigure{figures/Symbolic Axis Mapping.pdf}{Combining local axis mappings for dot and addition. The shared tensor $T$ establishes $\vec S=\vec I\catseq\vec J$, allowing the compiler to propagate the positions of $\vec I$ and $\vec J$ to $Z$ and $O$.}{fig:axis-mapping}{Local dot and addition mapping tables are connected through their shared
intermediate, producing a table whose axes use a common coordinate description.}
The core question is: how can we relate different symbolic axis groups
across composed operations? The compiler must track the same logical axes
across operations and infer their positions in each tensor. These positions
are needed to construct RHS attributes: for example, contracting a group
requires its positions in both operands.

Each mapping stores a triple $(\mathsf{I}_t,\vec G,q)$: $\mathsf{I}_t$ is
the ordered list of logical axes in tensor $t$, $\vec G$ is the symbolic group
being located, and $q$ lists the positions of its axes in $\mathsf{I}_t$,
in group order. Thus
$\gath(\mathsf{I}_t,q)=\vec G$ (Equation~\eqref{eq:mapping}).
The positions may be symbolic expressions rather than concrete lists.

We write $\rho(\vec G)=\operatorname{range}(|\vec G|)=[0,1,\ldots,|\vec G|-1]$ for
the ordered list of positions within $\vec G$, where $|\vec G|$ is its number of axes.
For example, if $\vec G=[i_0,i_1,i_2]$, then $\rho(\vec G)=[0,1,2]$.
These are positions within the group; the mapping list $q$ gives their
positions in tensor $t$. For an empty group, $\rho(\vec G)=[]$.

The compiler builds mappings from producers to consumers in three steps:
\begin{enumerate}
\item \emph{Record local mappings.}
Use each operation's \code{bind} declarations and output order to locate
its groups in its operands and result.
\item \emph{Connect shared interfaces.}
Align the producer's output groups with the consumer's description of
the same tensor, preserving axis order and pairings.
\item \emph{Compose positions.}
An operation may select axes from a group whose positions in the original
tensor are already known. To find those axes in the tensor, the compiler
combines the two position lists. For example, if $\vec G$ occupies tensor
positions $q=[2,0,1]$ and the operation selects positions $h=[1,0]$
within $\vec G$, their tensor positions are $\gath(q,h)=[0,2]$.
In general,
\begin{equation}
 \gath\!\left(
   \underbrace{\gath(\mathsf{I}_t,q)}_{\text{group }\vec G},
   \underbrace{h}_{\text{positions within }\vec G}
 \right)
 =
 \gath\!\left(
   \mathsf{I}_t,
   \underbrace{\gath(q,h)}_{\text{positions in }\mathsf{I}_t}
 \right).
 \label{eq:gather-composition}
\end{equation}
This means that the compiler can combine position mappings across composed
operations to locate selected axes in the corresponding tensors.
\end{enumerate}

Figure~\ref{fig:axis-mapping} combines the local mappings of dot and
addition through their shared tensor $T$. Identifying $\vec S=\vec I\catseq\vec J$
lets the compiler infer the positions of $\vec I$ and $\vec J$ in $Z$ and $O$.

\subsubsection{Symbolic Axis Partitioning}
\label{sec:partitioning}
A consumer may select axes from several producer groups. How can the
compiler separate these selections while preserving their order and operand
pairings? Symbolic axis partitioning divides each producer group into
selected and remaining axes and records their positions.

For dot reassociation $(XY)Z\to X(YZ)$, the inner dot contracts $\vec K$ and
produces $T$ with groups $\vec B\catseq\vec I\catseq\vec J$. The outer dot may contract
axes from both $\vec I$ and $\vec J$. Figure~\ref{fig:axis-partitioning} selects
$[j_1,i_0]$, paired with positions $[2,0]$ in $Z$. The compiler separates
these selections by their producer groups while retaining each axis's
partner in $Z$. We assume that the \code{where} block enforces a guard
requiring the outer dot's left batch list to select the inner batch group
in the same order.

\paperfigure{1}{figures/Symbolic Axis Partitioning (Dot Reassociation).pdf}{Partitioning selected axes by producer origin preserves their contraction partners for reassociation.}{fig:axis-partitioning}{Three panels show selecting producer output axes, splitting them by origin,
and retaining contraction partners, followed by a subgroup mapping table.}
Let $k_\ell^{\mathrm{outer}}$ and $k_r^{\mathrm{outer}}$ be the outer dot's
contraction lists in $T$ and $Z$, pairing axes of the contraction group $\vec D$.
For each producer group $\vec G\in\{\vec I,\vec J\}$, the compiler records three lists:
\begin{itemize}
\item $v_G$: positions of selected axes within $\vec G$;
\item $u_G$: positions of those axes in the outer contraction list;
\item $z_G$: positions of their paired axes in $Z$.
\end{itemize}
The operation $\mpo(a,b)=(u,v)$ finds the entries shared by
distinct-entry lists $a$ and $b$, in $a$'s order, and returns their positions
in each list. Thus $\gath(a,u)=\gath(b,v)$.
Let $q_G$ locate $\vec G$ in $T$. The compiler computes
\begin{equation}
\begin{aligned}
 (u_G,v_G)&=\mpo(k_\ell^{\mathrm{outer}},q_G),\\
 \vec G_{\mathrm{sel}}&=\gath(\vec G,v_G),\qquad
 \vec G_{\mathrm{rem}}=\gath(\vec G,\ocomp(\rho(\vec G),v_G)),\\
 z_G&=\gath(k_r^{\mathrm{outer}},u_G).
\end{aligned}
\label{eq:z-partners}
\end{equation}
Selected axes follow contraction order, while remaining axes preserve
their order in $\vec G$. Concatenating the two subgroups therefore need not
restore $\vec G$'s original order.

For $i_0$ in the figure, its position within $\vec I$ is 0, but its position
in the contraction list is 1. Thus $v_I=[0]$ and $u_I=[1]$.
Its partner is found at entry 1 of the right contraction list $[2,0]$,
giving $z_I=[0]$. Using $v_I$ instead of $u_I$ would select the wrong partner.

These refined mappings let the compiler locate each subgroup in the
replacement operands and construct attributes that preserve the original
axis pairings, even when group lengths and positions remain symbolic.
Appendix~\ref{extended:sec:partitioning} gives the detailed batch condition,
output order, and complete refined mapping table.

\subsubsection{Symbolic Axis Mapping Algorithm}
\label{sec:mapping-algorithm}
The compiler builds shared axis mappings by visiting producers before
consumers. At each operation, it connects the operation's local axis
bindings to the mappings already established for its inputs.

The mapping procedure performs four steps:
\begin{enumerate}
\item Read the operation's axis bindings and locate each group in its operands.
\item Connect these groups to the input mappings, splitting groups when
      the producer and consumer describe the same axes differently.
\item Express operand reads using the aligned groups and combine them
      according to the operation's value expression.
\item Record the output axis order and the positions of its groups.
\end{enumerate}

The compiler keeps each tensor's axis interface and value expression in
\code{V}. It also records mappings that locate individual groups and
partitions that describe how groups together cover an entire interface.
For example, an interface may be described by one group $\vec S$ or by two
groups $\vec I$ and $\vec J$, where $\vec S=\vec I\catseq\vec J$. The compiler connects these
descriptions while preserving the tensor's axis order.

The resulting mappings let attribute synthesis locate logical axes in
replacement operands. Appendix~\ref{extended:sec:mapping-algorithm}
presents the full algorithm and its construction invariants.
